\documentclass[10pt, aspectratio=1610]{beamer}
\usetheme{UCL}
\usepackage{bookmark}

\hypersetup{
  colorlinks=true,
  linkcolor={uclgold},
  filecolor={Maroon},
  citecolor={uclgold},
  urlcolor={uclgold}
}

%------------------------------------------------------------
% PRESENTATION METADATA
%------------------------------------------------------------
\title[A short title]{A very long title for my slides, which may span across two lines}
\author{Your Name}
\date[20 Jul 2026]{20 July 2026}

% Using your new automated conference metadata macro
\conference[SHORTER TITLE OF THE CONFERENCE]{The full name of the conference, Where it is\\ The name of the session}

% Affiliation
\institute{Department of Statistical Science ~ | ~ University College London}

% These can be commented out if you don't want email/website/social media links
\newcommand{\insertsocials}{%
  {\setlength{\baselineskip}{9pt}%
  \faEnvelope~\href{mailto:n.surname@ucl.ac.uk}{\texttt{n.surname@ucl.ac.uk}} \\
  \faGlobe~\href{https://yourwebsite.com}{\texttt{https://yourwebsite.com}}\\
  \faGithub~\href{https://github.com/you}{\texttt{https://github.com/you}} \\
% Can also add social media links, e.g.
  \faMastodon~\href{https://mas.to/@yourmastodon}{\texttt{@yourmastodon@mas.to}} \,\textbar\, 
  \faLinkedinIn~\href{https://www.linkedin.com/in/name-surname/}{\texttt{name-surname}}
  }%
}


%------------------------------------------------------------
% PRESENTATION BODY
%------------------------------------------------------------
\begin{document}

\frame[plain]{\titlepage}

\begin{frame}{First Slide Title}
  The footer below now displays text-only metadata flawlessly.
\end{frame}

\begin{frame}{Hello there {\textemoji{👋}}}
\begin{itemize}
\item
  Presenting code and LaTeX equations
\item
  Images
\end{itemize}

\ldots and much more
\end{frame}

\begin{frame}{No logo}
  \global\uclNoLogotrue
  Here's a slide with no logo -- the footline macro safely cleans this frame.
\end{frame}

\begin{frame}[fragile]{Some maths}
  We can also include maths, using standard \LaTeX:

  \[\class{italian-green}{p(y\mid\mathcal{D}) = \int p(y\mid\bm\theta)p(\boldsymbol\theta\mid\mathcal{D})d\boldsymbol\theta} \]

  To use colours with \LaTeX, you can utilize the custom theme wrapper:
  \texttt{\$\textbackslash{}class\{color\}\{expression\}\$}

  You can define a probabilistic assumption as
  \[\class{spanish-red}{y \mid \bm\theta \sim \dpois(\bm\theta)}\] 
  but in line it would look like this \(\mathcolor{myblue}{y\sim \dnorm(\mu,\sigma^2)}\).
\end{frame}

\section{Example Section Divider}

\begin{frame}{After Section Frame}
  This slide appears directly after a section break layout has rendered automatically.
\end{frame}

\begin{frame}[fragile,noframenumbering]{}
\uclDividerFrameStart{A section}

Using the tag \verb"\uclDividerFrameStart" makes for some nicer formatting (must be turned off before the end of the frame with
\verb"\uclDividerFrameEnd").

\uclDividerFrameEnd
\end{frame}

\begin{frame}{Column Layout}
\protect\phantomsection\label{column-layout}
Arrange content into columns of varying widths:

\begin{columns}[T]
\begin{column}{0.35\linewidth}
\begin{block}{Motor Trend Car Road Tests}
\protect\phantomsection\label{motor-trend-car-road-tests}
The data was extracted from the 1974 Motor Trend US magazine, and
comprises fuel consumption and 10 aspects of automobile design and
performance for 32 automobiles.
\end{block}
\end{column}

\begin{column}{0.03\linewidth}
\end{column}

\begin{column}{0.62\linewidth}
{\def\LTcaptype{none} % do not increment counter
\begin{longtable}[]{@{}lrrrrr@{}}
\toprule\noalign{}
& mpg & cyl & disp & hp & wt \\
\midrule\noalign{}
\endhead
Mazda RX4 & 21.0 & 6 & 160 & 110 & 2.620 \\
Mazda RX4 Wag & 21.0 & 6 & 160 & 110 & 2.875 \\
Datsun 710 & 22.8 & 4 & 108 & 93 & 2.320 \\
Hornet 4 Drive & 21.4 & 6 & 258 & 110 & 3.215 \\
Hornet Sportabout & 18.7 & 8 & 360 & 175 & 3.440 \\
Valiant & 18.1 & 6 & 225 & 105 & 3.460 \\
\bottomrule\noalign{}
\end{longtable}
}
\end{column}
\end{columns}

Learn more:
\href{https://quarto.org/docs/presentations/revealjs/\#multiple-columns}{Multiple
Columns}
\end{frame}

\end{document}
